\section{A1: Identical bosons and symmetric wavefunctions}

\subsection{Exchange of identical particles}

For two distinguishable particles in one-particle states
$\phi_a$ and $\phi_b$, a product state is
\[
\Psi(\mathbf r_1,\mathbf r_2)
=\phi_a(\mathbf r_1)\phi_b(\mathbf r_2).
\]
For identical particles, the labels $1$ and $2$ label coordinate slots rather
than physically distinguishable atoms. Define the exchange operator by
\[
\hat P_{12}\Psi(\mathbf r_1,\mathbf r_2)
=\Psi(\mathbf r_2,\mathbf r_1).
\]
Since two exchanges restore the original state,
\[
\hat P_{12}^2=\hat I.
\]
If $\hat P_{12}\Psi=\lambda\Psi$, then
$\lambda^2=1$, so $\lambda=\pm1$. Bosons occupy the symmetric sector:
\[
\boxed{\hat P_{12}\Psi_B=+\Psi_B.}
\]

\subsection{Two bosons in different orbitals}

For normalized orthogonal orbitals $a$ and $b$, the symmetric state is
\[
\boxed{
\Psi_B(\mathbf r_1,\mathbf r_2)
=\frac{1}{\sqrt2}
\left[
\phi_a(\mathbf r_1)\phi_b(\mathbf r_2)
+\phi_b(\mathbf r_1)\phi_a(\mathbf r_2)
\right].
}
\]
The normalization follows from
\[
\langle\Psi_B|\Psi_B\rangle
=|C|^2(1+0+0+1)=2|C|^2=1.
\]
The state says that one boson occupies $a$ and one occupies $b$; there is no
physical fact assigning a particular atom to either orbital.

Squaring the wavefunction produces an exchange-interference term:
\[
|\Psi_B|^2
=\frac12|a_1b_2|^2+\frac12|b_1a_2|^2
+\operatorname{Re}(a_1^*b_2^*b_1a_2).
\]
Thus indistinguishability affects measurable probabilities even without a
dynamical interaction between the particles.

\subsection{Two bosons in the same orbital}

If both particles occupy $\phi_a$, the normalized wavefunction is
\[
\boxed{
\Psi_B(\mathbf r_1,\mathbf r_2)
=\phi_a(\mathbf r_1)\phi_a(\mathbf r_2).
}
\]
For $N$ bosons in one orbital $\phi_0$,
\[
\boxed{
\Psi_N(\mathbf r_1,\ldots,\mathbf r_N)
=\prod_{j=1}^{N}\phi_0(\mathbf r_j).
}
\]
This is automatically symmetric. The ability of arbitrarily many bosons to
occupy one orbital is the statistical foundation of Bose--Einstein
condensation.

\subsection{Symmetrization operator}

For two particles,
\[
\hat S=\frac12(\hat I+\hat P_{12}).
\]
It satisfies
\[
\hat S^2=\hat S,
\qquad
\hat P_{12}\hat S=\hat S,
\]
so it projects a two-particle state onto its symmetric part. Projection and
normalization are separate operations.

\paragraph{Checkpoint.}
The labels $1,2$ identify wavefunction arguments, not individual bosons. The
natural physical information is therefore the number of bosons occupying each
single-particle orbital.
